\chapter{Remote HPC Execution}
\label{ch:remote}

The \ui{HPC} panel (zone 11) submits calculations to a cluster
over SSH, monitors the queue, and brings the results back --- without
leaving Calango.

\begin{calwarn}
Needs \texttt{paramiko} in Calango's Python environment, and SSH access to
a cluster running SLURM, PBS or SGE.
\end{calwarn}

\begin{calnote}
All network traffic runs in a separate helper process driven over a JSON
protocol, mirroring how local jobs are isolated: the interface never blocks
on the network, and a hung connection can always be resolved by killing
the helper rather than the application.
\end{calnote}

\begin{calnote}
The connection is a \emph{session}, not a login per operation: Calango
authenticates once and multiplexes every upload, submission, status poll,
log tail and download over that one SSH transport. On a cluster with
keyboard-interactive two-factor authentication this is the difference
between usable and unusable --- the one-time code is typed once, not once
per queue poll.
\end{calnote}

\section{Connection}

The \ui{Connection} tab holds the cluster credentials:

\begin{description}
  \item[\ui{Host}, \ui{Port}] Hostname and port (default 22).
  \item[\ui{User}] Your account name on the cluster.
  \item[\ui{Auth}] \ui{SSH key} or \ui{Password}.
  \item[\ui{Key file}] Path to a private key, e.g.\
    \file{\textasciitilde/.ssh/id\_rsa}. Leave empty to let your SSH agent
    or default keys authenticate.
  \item[\ui{Password}] The password, or the passphrase for an encrypted
    key.
  \item[\ui{Remote dir}] Working directory on the cluster, relative to
    \opt{\$HOME} (default \file{calango\_jobs}).
  \item[\ui{Connect}] Opens the session, reports your remote
    \opt{\$HOME}, and auto-detects the scheduler. Pressing it first is
    optional: any operation issued while disconnected authenticates and
    then runs.
  \item[\ui{Disconnect}] Closes the session. A job already submitted keeps
    running on the cluster.
\end{description}

\subsection*{Two-factor authentication}

If the cluster asks for a second factor, Calango puts the server's own
challenge on screen --- its wording, its prompts, with the answer masked
whenever the server says it must not be echoed --- and sends what you type
straight back. Clusters that ask for the password and the code in one
exchange get the password filled in from the \ui{Password} field, so only
the code is asked for.

The code is never stored, which is also why a session opened with 2FA is
never re-established silently: if it drops, the panel says so and waits for
you to press \ui{Connect}. A session opened with a key or a password has
nothing single-use about it and reconnects by itself on the next operation.

\begin{calwarn}
Passwords and key passphrases are kept in memory only --- never written to
disk, never placed on a command line where \texttt{ps} could see them.
One-time codes are not kept even that long. Everything else (host, user,
key path, scheduler settings) is remembered between sessions. A host key
that \emph{changes} is refused outright rather than pinned over.
\end{calwarn}

\section{Scheduler settings}

The \ui{Scheduler} tab describes the batch job:

\begin{description}
  \item[\ui{Scheduler}] \ui{SLURM}, \ui{PBS} or \ui{SGE}.
  \item[\ui{Queue}] Partition or queue name; empty uses the cluster default.
  \item[\ui{Tasks / cores}] Requested cores.
  \item[\ui{Walltime}] \texttt{HH:MM:SS}, default \texttt{01:00:00}.
  \item[\ui{Setup}] Shell lines executed before the payload --- module
    loads, \texttt{conda activate}, anything your site needs:
    \begin{lstlisting}
module load python
source ~/venvs/ase/bin/activate
    \end{lstlisting}
  \item[\ui{Submit Calculation\ldots}] Opens the usual calculator dialog,
    then submits the result.
\end{description}

The generated wrapper carries the right directives for your scheduler ---
\texttt{\#SBATCH}, \texttt{\#PBS} or \texttt{\#\$} --- and always redirects
output to fixed file names so the monitor knows what to tail.

\section{Submitting and monitoring}

\menu{Simulation}\menusep\menu{New Remote Calculation}
(\kbd{Ctrl}+\kbd{Shift}+\kbd{R}) --- or \ui{Submit Calculation\ldots} in
the panel --- runs the full sequence:

\begin{enumerate}
  \item Stage \file{run.py} and \file{structure.extxyz} into a local job
    directory, exactly as a local run would.
  \item Generate \file{job.sh} from the scheduler settings.
  \item Upload the directory over SFTP, creating the remote path as needed.
  \item Submit with \texttt{sbatch} or \texttt{qsub} and capture the job ID.
  \item Start monitoring.
\end{enumerate}

The \ui{Queue \& Logs} tab then shows the job ID and remote directory, the
live queue state polled with \texttt{squeue} or \texttt{qstat}, and the
remote standard output and error streamed incrementally into the console
(errors in red). \ui{Cancel Job} issues \texttt{scancel} or \texttt{qdel}.

\section{Retrieving results}

When the job leaves the queue, Calango downloads the results automatically
--- structures, trajectories, logs and CSV metric files --- into the same
local job directory the inputs came from. Trajectories and final geometries
open directly in a new tab.

\ui{Download Results} repeats the transfer manually, which is useful for
pulling intermediate output from a long-running job without waiting for it
to finish.

The task appears in the process manager like any other, so its directory
stays one click away for later analysis.
